\documentclass[11pt]{article}
\usepackage{../../estilo/vanguarda44}

\renewcommand{\contentsname}{Sumário}

\title{\bfseries\color{vinho}\Huge VANGUARDA 44: NORMANDIA}
\author{}
\date{}

\hypersetup{
  pdftitle={Vanguarda 44: Normandia -- Mini-Livro de Expansão},
  pdfauthor={Felipe Ramires Terrazas},
}

\fancyhead[L]{\small\color{cinzachumbo}Vanguarda 44: Normandia --- Mini-Livro de Expansão}

\begin{document}

\begin{titlepage}
\centering
\vspace*{3cm}
{\Huge\bfseries\color{vinho} VANGUARDA 44: NORMANDIA \par}
\vspace{0.5cm}
{\Large\color{mostarda} Mini-Livro de Expansão \par}
\vspace{2cm}
{\large Doutrina Britânica e Cenários Históricos da Campanha da Normandia \par}
\vfill
\end{titlepage}

\tableofcontents

\bigskip

Junho de 1944. As areias de Omaha, os campos alagados ao redor de Carentan, as ruínas bombardeadas de Caen --- cada palmo da Normandia custou caro aos dois lados desta campanha que abriu a Frente Ocidental. Este mini-livro traz a doutrina britânica, ausente do Livro de Regras principal, e três cenários históricos para você revivê-la à mesa.

\section{Doutrina Nacional: Reino Unido}

O exército britânico construiu sua doutrina de esquadrão ao redor da Bren, uma metralhadora leve confiável e portátil o suficiente para acompanhar o avanço da infantaria --- diferente da doutrina alemã, que mantém a MG42 estática, ou da americana, que aposta no fogo individual superior do fuzileiro.

\textbf{Regra do Exército}: a Metralhadora Leve britânica (Bren) NÃO sofre a penalidade de -1 no Teste de Acerto ao mover com a unidade --- ela é parte integral do avanço do esquadrão, não uma arma de apoio estático.

\subsection{Sub-Facção A: Infantaria de Linha}

Foco: disciplina defensiva e cobertura mútua. Regra Especial (``Mantenham a Posição''): esquadrões britânicos que não se moveram neste turno ignoram o primeiro nível de Stress recebido (empilha com o bônus de Veterano, se aplicável).

\subsection{Sub-Facção B: Comandos / Pára-quedistas Britânicos}

Foco: ataques relâmpago coordenados, como a captura da Pegasus Bridge nas primeiras horas do Dia-D. Regra Especial (``Ataque Relâmpago''): no primeiro turno da partida, esses esquadrões executam Avançar sem a penalidade normal de -1.

\section{Cenários}

\subsection{Dia-D --- Praia de Omaha}

\textbf{Contexto Histórico}: Na madrugada de 6 de junho de 1944, milhares de soldados americanos desceram de suas embarcações de desembarque direto numa praia aberta, sem cobertura alguma, contra a 352ª Divisão de Infantaria alemã entrincheirada no alto dos penhascos. O Alto Comando alemão permanecia convencido de que o ataque principal viria em Pas-de-Calais, e essa incerteza atrasou a resposta organizada da defesa nas horas mais críticas do desembarque --- uma janela de confusão que os Aliados exploraram a cada metro conquistado na areia.

\textbf{Data e Local}: 6 de junho de 1944, Praia de Omaha, Normandia, França.

\textbf{Forças}: Americanos (Atacante) vs. Alemães (Defensor).

\textbf{Terreno}: uma faixa de Campo Aberto (a praia) ocupando o terço da mesa mais próximo da implantação americana; o restante com Obstáculos Lineares (obstáculos de praia) e uma linha de Bunker/Fortificação no lado alemão, no topo dos penhascos.

\textbf{Implantação}: Americanos entram pela borda da mesa mais próxima da água (desembarcando), sem posição fixa; Alemães começam já posicionados nos bunkers e trincheiras.

\textbf{Qualidade de Tropa}: Alemães Regulares (guarnição costeira); Americanos Regulares.

\textbf{Regra Especial --- ``A Surpresa da Invasão''}: nos primeiros 2 turnos, unidades americanas que executarem Avançar ou Correr recebem +5cm extra de movimento, representando o avanço através da praia antes que a defesa alemã se reorganize por completo.

\textbf{Objetivo}: Americanos vencem se ocuparem a linha de bunkers até o fim do turno 8. Alemães vencem se os impedirem.

\textbf{Duração}: 8 turnos.

\subsection{Carentan}

\textbf{Contexto Histórico}: A 101ª Divisão Aerotransportada avançou sobre Carentan --- ponto crucial para conectar as praias de Utah e Omaha --- através de uma estrada estreita cercada de campos alagados que os próprios soldados apelidaram de ``Purple Heart Lane''. A cidade era defendida pelo Regimento de Paraquedistas 6, uma unidade de elite alemã sob o comando do Coronel Friedrich von der Heydte, veterano de Creta. Foram quatro dias de combate quase ininterrupto pela estrada e pelas ruas da cidade antes que os americanos a tomassem definitivamente.

\textbf{Data e Local}: 10 a 13 de junho de 1944, Carentan, Normandia, França.

\textbf{Forças}: Americanos (Atacante) vs. Alemães (Defensor).

\textbf{Terreno}: uma estrada/causeway estreita (cerca de 2cm de largura na mesa) cercada de Terreno Difícil ou área intransponível (os campos alagados) dos dois lados, conectando a implantação americana à cidade; a cidade em si com Cobertura Pesada.

\textbf{Implantação}: Americanos numa borda estreita da mesa (a entrada da estrada); Alemães na cidade.

\textbf{Qualidade de Tropa}: ambos os lados Veteranos (paraquedistas de elite dos dois exércitos).

\textbf{Regra Especial --- ``Em Guarda''}: unidades alemãs que dispararem contra um alvo americano que executou Avançar ou Correr neste turno ganham +1 no Teste de Acerto --- os defensores já esperavam o ataque pela estrada, com armas pré-posicionadas e comunicação estabelecida.

\textbf{Objetivo}: Americanos vencem se ocuparem o centro da cidade até o turno 8. Alemães vencem se resistirem.

\textbf{Duração}: 8 turnos.

\subsection{Caen}

\textbf{Contexto Histórico}: Caen era um objetivo previsto para o próprio Dia-D, mas só caiu por completo mais de um mês depois, num combate urbano prolongado e destrutivo. A Operação Charnwood, em julho de 1944, viu forças britânicas e canadenses avançarem casa a casa contra uma defesa alemã tenaz, apoiada por blindados, através de uma cidade já reduzida a escombros por semanas de bombardeio aéreo.

\textbf{Data e Local}: julho de 1944, Caen, Normandia, França.

\textbf{Forças}: Britânicos (Atacante) vs. Alemães (Defensor).

\textbf{Terreno}: predominantemente urbano --- Cobertura Pesada em quase toda a mesa, com ruínas de bombardeio.

\textbf{Implantação}: Britânicos avançam de um lado da cidade; Alemães defendem em profundidade, com múltiplas linhas de defesa (não uma única linha de frente).

\textbf{Qualidade de Tropa}: Alemães Veteranos (unidades blindadas endurecidas em combate); Britânicos Regulares.

\textbf{Regra Especial --- ``Bombardeio Prévio''}: no início da partida, role 1D6 para cada peça de terreno de edifício na mesa; num resultado de 5+, aquele edifício já está em ruínas --- mantém Cobertura Pesada, mas também vira Terreno Difícil para atravessar.

\textbf{Objetivo}: Britânicos vencem se controlarem a maioria de 4 a 6 zonas de objetivo definidas na mesa até o turno 10.

\textbf{Duração}: 10 turnos.

\end{document}
